% Exploratory cross-study analyses: where each one is reported, and the analyses that no
% Supplementary Note covers. Numbers come only from macros in generated/numbers.tex.
\label{sec:suppB-crossstudy}

Every analysis in this section is exploratory. None is named in any registration: the pooled analysis of penalties and the two analyses of explanations were specified on 9 October 2026, and the image check on 10 October 2026, after every registered result was known. Their tests are not corrected for multiple comparisons, and they change no registered verdict. The pooled analyses treat each study's participants as different people, so a person who took part in Study~5 and in Study~2 or~4 counts as two (Section~\ref{sec:suppB-study5}). A $p$-value is the chance of a result at least this large if there were no effect.

Most of these analyses are reported elsewhere. The check of the stimulus images is in Supplementary Note~4. Supplementary Note~6 and Supplementary Fig.~\ref{figS:pooled} give the model that pools the label's penalty on second works over Studies 2, 4 and 5, with the comparison of rival accounts and the results when one study is left out. Supplementary Note~7 gives the analyses of what the explanations referred to, including how often they mentioned prompting or a machine and how far they moved from None to Dramatic. Supplementary Note~5, Supplementary Fig.~\ref{figS:editingtime} and Supplementary Table~\ref{tab:suppB-s5breakdown} give Study~5's penalty at each stated editing time. This section reports the rest: the shape of the penalty within each study, a check of the pooled model, Study~4's penalty at each stated time, the explanation shift on the two measures other than creativity, and a structural topic model of the explanations.

\paragraph*{The shape of the penalty within each study.}
Here the penalty is the predicted mean rating of a second work under the Human label minus the predicted mean under the AI label. Supplementary Table~\ref{tab:suppB-xstrend} gives its slope and curvature across the three levels in each study, from the pooled model. The slope is the change in penalty per step from None to Dramatic, which dilution, the account on which the penalty is largest at None and shrinks as the edit grows, predicts to be negative. The curvature measures how far the Slight penalty stands above a straight line through the None and Dramatic penalties. In no study did the penalty fall significantly from None to Dramatic on any measure. It rose significantly before correction in two places: Study~2's creativity penalty (slope \XSCreatSlopeSTwoEst{} points per step, $p$~\XSCreatSlopeSTwoP) and Study~5's transformativeness penalty (\XSTransSlopeSFiveEst, $p$~\XSTransSlopeSFiveP). The Study~2 creativity slope is Study~2's registered H3 contrast with its sign reversed and halved, refitted in a model that gives each second work its own baseline; the registered test, which was not significant (Holm $p$~\STwoCreatPenaltyChangePHolm; Supplementary Note~6), governs. The curvature was positive, a peak at Slight, and clear of zero before correction on essence change in all three studies, on transformativeness in Study~4 only, and on creativity in Study~5 only, with Study~4 close. Of 24 uncorrected tests, about one in twenty can fall below .05 by chance alone.

% Exploratory: slope and curvature of the label's penalty across levels, in each study and pooled,
% from the pooled model of Studies 2, 4 and 5.
% Every number is a macro from generated/numbers.tex; supplement/tables/README_suppB.md lists
% the pipeline tables behind them.
\begin{table}[htbp]
\centering
\footnotesize
\caption{\textbf{Exploratory: the shape of the label's penalty across levels, study by study.} The penalty is the predicted mean rating under the Human label minus the predicted mean under the AI label. The slope is the change in penalty per step from None to Dramatic, half of the Dramatic penalty minus the None penalty; dilution predicts it to be negative. The curvature is the Slight penalty minus the average of the None and Dramatic penalties; it is zero for any straight line, and positive when the penalty peaks at Slight. All values come from the pooled model of Studies 2, 4 and 5, specified on 9 October 2026 after the registered results were known. The $p$-values are two-sided and not corrected for the 24 tests.}\label{tab:suppB-xstrend}
\setlength{\tabcolsep}{8pt}
\begin{tabular}{@{}lcccc@{}}
\toprule
 & \multicolumn{2}{c}{Slope} & \multicolumn{2}{c}{Curvature}\\
\cmidrule(lr){2-3}\cmidrule(l){4-5}
Scope & Estimate & Test & Estimate & Test\\
\midrule
\multicolumn{5}{@{}l}{\textit{Transformativeness}}\\
Study~2 & \XSTransSlopeSTwoEst & $p$~\XSTransSlopeSTwoP & \XSTransCurvatureSTwoEst & $p$~\XSTransCurvatureSTwoP\\
Study~4 & \XSTransSlopeSFourEst & $p$~\XSTransSlopeSFourP & \XSTransCurvatureSFourEst & $p$~\XSTransCurvatureSFourP\\
Study~5 & \XSTransSlopeSFiveEst & $p$~\XSTransSlopeSFiveP & \XSTransCurvatureSFiveEst & $p$~\XSTransCurvatureSFiveP\\
Pooled & \XSTransSlopeEst & $p$~\XSTransSlopeP & \XSTransCurvatureEst & $p$~\XSTransCurvatureP\\
\addlinespace[3pt]
\multicolumn{5}{@{}l}{\textit{Essence change}}\\
Study~2 & \XSEssSlopeSTwoEst & $p$~\XSEssSlopeSTwoP & \XSEssCurvatureSTwoEst & $p$~\XSEssCurvatureSTwoP\\
Study~4 & \XSEssSlopeSFourEst & $p$~\XSEssSlopeSFourP & \XSEssCurvatureSFourEst & $p$~\XSEssCurvatureSFourP\\
Study~5 & \XSEssSlopeSFiveEst & $p$~\XSEssSlopeSFiveP & \XSEssCurvatureSFiveEst & $p$~\XSEssCurvatureSFiveP\\
Pooled & \XSEssSlopeEst & $p$~\XSEssSlopeP & \XSEssCurvatureEst & $p$~\XSEssCurvatureP\\
\addlinespace[3pt]
\multicolumn{5}{@{}l}{\textit{Creativity}}\\
Study~2 & \XSCreatSlopeSTwoEst & $p$~\XSCreatSlopeSTwoP & \XSCreatCurvatureSTwoEst & $p$~\XSCreatCurvatureSTwoP\\
Study~4 & \XSCreatSlopeSFourEst & $p$~\XSCreatSlopeSFourP & \XSCreatCurvatureSFourEst & $p$~\XSCreatCurvatureSFourP\\
Study~5 & \XSCreatSlopeSFiveEst & $p$~\XSCreatSlopeSFiveP & \XSCreatCurvatureSFiveEst & $p$~\XSCreatCurvatureSFiveP\\
Pooled & \XSCreatSlopeEst & $p$~\XSCreatSlopeP & \XSCreatCurvatureEst & $p$~\XSCreatCurvatureP\\
\bottomrule
\end{tabular}

\smallskip
\begin{minipage}{\textwidth}
\footnotesize\raggedright
The pooled model gives each participant, each original and each second work its own baseline, and it leaves out Study~4's stated hours for the original and Study~5's stated editing hours, so each study's values average over them. Its fits converged but are singular: one source of variation, the spread between originals once each second work has its own baseline, was estimated at zero. Supplementary Note~6 gives the pooled slopes and curvatures with their 95\% confidence intervals.
\end{minipage}
\end{table}


\paragraph*{A check of the pooled model.}
The pooled model leaves out the time factors of Studies 4 and 5 and adds a baseline for each second work, so its estimate of each study's penalty need not equal that study's registered estimate. A descriptive check compared them. At every level and on every measure, the pooled model's estimate lay within a few thousandths of a point of the registered estimate. Only the point estimates agree this closely; the pooled model's standard errors, and so its $p$-values, differ from each study's own.

\paragraph*{Study~4's penalty at each stated time.}
A refit of Study~4's registered model estimated the label's penalty at each of the three hours stated for the original and each level, on each measure, without correction. It showed no consistent pattern across the stated hours. Of its \XSTimeCheckSFourNCells{} cells, \XSTimeCheckSFourNExcludeZero{} had a 95\% confidence interval that excluded zero: essence change at Dramatic for 100-hour originals, where second works from AI-labeled originals were rated higher (penalty \XSTimeCheckSFourEssPenaltyHundredHoursDramaticEst{} points; 95\% confidence interval \XSTimeCheckSFourEssPenaltyHundredHoursDramaticLo{} to \XSTimeCheckSFourEssPenaltyHundredHoursDramaticHi). With this many uncorrected intervals, a cell like this can arise by chance. None of Study~4's registered tests of the stated hours was significant (Supplementary Table~\ref{tab:suppB-s4}b). The cell estimates are in the cross-study tables that the analysis code writes (main text, Code availability).

\paragraph*{How far explanations moved, on transformativeness and essence change.}
Under taint, the account on which an AI origin passes to anything built from it, AI-label raters should keep writing about the original, so their explanations should change less from None to Dramatic. Supplementary Note~7 describes the measure and gives the creativity estimates. On the other two measures, too, no difference between the label groups was shown. The pooled difference, AI minus Human, was \XSTransShiftEmbEst{} on transformativeness with sentence embeddings (95\% confidence interval \XSTransShiftEmbLo{} to \XSTransShiftEmbHi; permutation $p$~\XSTransShiftEmbPPerm) and \XSTransShiftWordEst{} with word overlap (\XSTransShiftWordLo{} to \XSTransShiftWordHi; $p$~\XSTransShiftWordPPerm). On essence change it was \XSEssShiftEmbEst{} with sentence embeddings (\XSEssShiftEmbLo{} to \XSEssShiftEmbHi; $p$~\XSEssShiftEmbPPerm) and \XSEssShiftWordEst{} with word overlap (\XSEssShiftWordLo{} to \XSEssShiftWordHi; $p$~\XSEssShiftWordPPerm). A sentence embedding turns an explanation into a list of numbers whose distances reflect differences in meaning; word overlap is the share of content words two explanations have in common. The permutation $p$-value is the share of label-shuffled refits, shuffled within each study, that gave a difference at least this large. The essence-change model with word overlap was singular, estimating one source of variation at zero, but it converged and keeps its numbers. On creativity, in both label groups, a person's Dramatic explanation lay closer to their own None explanation than other people's Dramatic explanations of the same image did, by similar amounts (descriptive). No difference shown is not the same as no difference.

\paragraph*{A structural topic model of the explanations.}
A topic model finds groups of words that tend to occur together, called topics, and describes each response as a mix of them. A structural topic model also lets the share of each topic depend on recorded features of the response, here the label, the level, their combination, the rating and the study. One model was fitted to each rating's explanations, pooled over Studies 2, 4 and 5. Words describing the task and each study's own item wording were removed, and only words used in at least five responses were kept. The number of topics was chosen between four and twelve by a rule fixed in code: the number whose model best predicted words set aside from fitting. No person or language model named or judged any topic; each is described by its most probable words.

The explanations were short, and responses that kept no word after this pruning were dropped. At None, the AI-label group lost a larger share of its responses this way than the Human-label group, so a label difference at None compares responses selected differently in the two groups.

On creativity, the topic whose most probable words were `ai', `similar', `feel' and `takes' took a larger share of AI-label explanations than of Human-label ones at None, by \XSTopicCreatSevenNoneEst{} (95\% confidence interval \XSTopicCreatSevenNoneLo{} to \XSTopicCreatSevenNoneHi), and at Slight, by \XSTopicCreatSevenSlightEst{} (\XSTopicCreatSevenSlightLo{} to \XSTopicCreatSevenSlightHi), but not at Dramatic (\XSTopicCreatSevenDramaticEst; \XSTopicCreatSevenDramaticLo{} to \XSTopicCreatSevenDramaticHi). A topic share runs from 0 to 1. The difference at Dramatic minus the difference at None was \XSTopicCreatSevenChangeEst{} (\XSTopicCreatSevenChangeLo{} to \XSTopicCreatSevenChangeHi). Three cautions apply. This topic was poorly recovered when the model was refitted from other starting points, so it may not survive another start. The intervals treat each person's many responses as independent, so they are likely too narrow. And the None comparison is affected by the uneven loss of responses described above. On transformativeness, at None, AI-label explanations gave more share than Human-label ones to a topic led by `look', `transformation' and `identical' and less to a topic led by `difference', `see' and `looks', and both differences shrank toward Dramatic. The model of the essence-change explanations did not converge within its limit of iterations, so it is not completed: its label effects are description only, with no interval, $p$-value or verdict.

Like every text analysis in the paper, these describe what people wrote after rating; they do not show why anyone rated as they did.
